\documentclass[12pt]{article}
\usepackage[utf8]{inputenc}
\usepackage[T1]{fontenc}
\usepackage{hyperref}
\title{Yoke Language Resources for K-12 Curriculum Development}
\author{Compiled by Claude Sonnet 5.5 (claude-sonnet-5-5) for Andr\'as Kornai}
\date{}
\begin{document}
\maketitle

\section{Introduction}
Yoke (ISO 639-3 yki; also called Yoki, Yauke, Pauwi, Jauke or Bitovondo) is a very poorly documented language of northern Papua, Indonesia. It is spoken in Mamberamo Hilir District, Mamberamo Raya Regency. Sources give roughly 200 speakers; Joshua Project gives 600. Classification varies: Glottolog treats it as Austronesian (Lower Mamberamo), while Donohue treats it as Papuan. The search found no educational materials.

\section{Primary Educational Resources}
None found. Joshua Project reports that Indonesian Malay and Yoke are both used at church and in the first year of schooling. This was not independently confirmed.

\section{Linguistic and Reference Materials}
\begin{itemize}
\item Yoke language (Wikipedia): \url{https://en.wikipedia.org/wiki/Yoke_language}
\item Lower Mamberamo languages (Wikipedia): \url{https://en.wikipedia.org/wiki/Lower_Mamberamo_languages}
\item Endangered Languages Project entry: \url{https://endangeredlanguages.com/node/87899}
\end{itemize}

\section{Religious and Community Resources}
\begin{itemize}
\item Joshua Project, Pauwi/Yoke people group (Bible translation listed as needed): \url{https://joshuaproject.net/people_groups/14346}
\item Joshua Project language page for yki: \url{https://joshuaproject.net/index.php/languages/yki}
\end{itemize}

\section{Digital and Audio Resources}
None found.

\section{Curriculum Development Context}
The language has no literacy tradition and is under pressure from the larger Warembori language. Curriculum work would need basic documentation first.
\end{document}